\subsection{Comfort performance across different resolution of \ac{gcm}}

\SetPicSubDir{ch4-Results/agentic/Result_Agentic}
\label{subsec:ComfPerf}

Figure~\ref{fig:ComfAgentic}\subref{fig:ComfAchi} shows the mean comfort probability under the \ac{sgcm}, \ac{dgcm}, \ac{rtgcm}, and fixed \qty{24}{\celsius} baseline policies for different subgroup sizes across the four office rooms.
Figure~\ref{fig:ComfAgentic}\subref{fig:ComfImp} shows the improvement in mean comfort probability relative to the fixed \qty{24}{\celsius} baseline.
Across all rooms, the baseline control remained nearly constant at approximately \qtyrange{65}{70}{\percent} across subgroup sizes.
The \ac{rtgcm} produced the highest mean comfort probability for every room and subgroup size, followed by the \ac{dgcm} and then the \ac{sgcm}.
The performance of the \ac{rtgcm} and \ac{dgcm} decreases at similar rates as subgroup size increases, while the performance of the \ac{sgcm} shows a more modest decrease.
The mean comfort probability across rooms ranged approximately \qtyrange{80}{98}{\percent} for the \ac{rtgcm}, \qtyrange{75}{95}{\percent} for the \ac{dgcm}, and \qtyrange{65}{73}{\percent} for the \ac{sgcm}.

The improvement relative to the 24 \si{\celsius} baseline decreased as subgroup size increased.
\ac{sgcm} improvement approaches zero in larger rooms, especially around subgroup size of 5--7 users.

In contrast, the \ac{rtgcm} and \ac{dgcm} maintain positive improvements at the largest tested subgroup size, with \ac{rtgcm} improvement remaining approximately \qtyrange{15}{20}{\percent} and \ac{dgcm} improvement approximately \qtyrange{5}{20}{\percent}.

Although both the \ac{dgcm} and \ac{rtgcm} maintain performance improvements over the baseline, the gap between them increases as room size expands.
This suggests that, in larger rooms, the \ac{rtgcm} benefits from more dynamic occupant transitions, whereas the \ac{dgcm} cannot fully capture time-varying group comfort as occupancy changes.

\begin{figure}[pos=htbp!]
    \centering
    \begin{subfigure}{0.98\linewidth}
        \centering
        \includegraphics[width=\linewidth]{\Pic{png}{a6p_static_realtime_daily_by_subgroup}}
        \caption{Mean comfort probability under different comfort-representation controls}
        \label{fig:ComfAchi}
    \end{subfigure}

    \medskip

    \begin{subfigure}{0.9\linewidth}
        \centering
        \includegraphics[width=\linewidth]{\Pic{png}{a6p_improvement_vs24_by_subgroup}}
        \caption{Comfort improvement relative to the baseline control}
        \label{fig:ComfImp}
    \end{subfigure}
    \caption{Agentic simulation results for different subgroup sizes.}
    \label{fig:ComfAgentic}
\end{figure}

% \autoref{fig:MeanOccuConf} shows the improvement ratio relative to the baseline control ($24^\circ\mathrm{C}$) as a function of daily mean occupancy. The figure presents the relationship between daily mean occupancy and improvement for the real-time and static group comfort models using density contours and fitted linear regression lines. Most observations are concentrated in regions where the improvement ratio is below 20\%. The fitted lines further indicate that the real-time group comfort model is more effective under sparse occupancy conditions. As mean occupancy increases to around 4--5 occupants, the improvement ratio approaches zero. This suggests that, even when room size (subgroup size) is fixed, the actual mean occupancy level still affects the achievable improvement.

% \begin{figure}[pos=htbp!]
%     \centering
%     \includegraphics[width=0.97\linewidth]{\Pic{png}{analysis72_mean_occupancy_contour_all_k4_k6_k8}}
%     \caption{Improvement ratio relative to the baseline control across daily mean occupancy levels}
%     \label{fig:MeanOccuConf}
% \end{figure}

\subsection{Effect of \ac{rtgcm} on control adjustment}
\label{subsec:controleffect}
The previous subsection showed that the \ac{rtgcm} maintained a \qtyrange{15}{30}{\percent} higher mean comfort probability than the \ac{sgcm} across room member sizes.
This subsection investigates the control-adjustment burden introduced by the \ac{rtgcm}.
Because the \ac{sgcm} and \ac{dgcm} use coarser occupant-composition representations, frequent setpoint updates are primarily required under \ac{rtgcm}-based control.

\autoref{fig:DynSpStepHist} shows the distribution of temperature adjustment magnitudes under the \ac{rtgcm} across daily mean occupancy.
As daily occupancy increases, the range of temperature adjustment magnitude slowly shrinks.
However, the mean value remains approximately around \qty{1.0}{\celsius} even at the highest occupancy levels.

Similarly, \autoref{fig:DynSpStepHeatmap} maps the temperature adjustment magnitude over subgroup size and mean occupancy ratio at the control timestep.
The magnitude of this temperature range shrinks as both subgroup size and mean occupancy ratio increase.
At small subgroup sizes such as $K=4$, the mean adjustment magnitude is over \qty{1.0}{\celsius} in most cases.


\begin{figure}[pos=htbp!]
    \centering
    \includegraphics[width=0.8\linewidth]{\Pic{png}{analysis23_mean_occupancy_moving_quantile_dynamic_sp_mean_abs_step_when_changed_expanded_pool_overlap_dense_only}}
    \caption{Distribution of temperature adjustment magnitudes under dynamic group control}
    \label{fig:DynSpStepHist}
\end{figure}

\begin{figure}[pos=htbp!]
    \centering
    \includegraphics[width=0.8\linewidth]{\Pic{png}{analysis23_k_normalized_mean_occupancy_mean_abs_step_when_changed_heatmap_regime_style_expanded_pool_overlap_dense_ybins05pct_kle16_minn20}}
    \caption{Mapping of temperature adjustment magnitude over subgroup size and mean occupancy ratio}
    \label{fig:DynSpStepHeatmap}
\end{figure}

\subsection{Effect of occupancy dynamics on control action frequency}

\autoref{fig:CtrlWindowRatio} shows the probability of temperature adjustment events of at least \qty{0.5}{\celsius} across different \ac{utr} levels.
As \ac{utr} increases, the probability of required temperature adjustment also increases; when \ac{utr} reaches approximately 0.7--0.8, the probability of temperature adjustment reaches \qty{90}{\percent}.

\autoref{fig:UTRexamples} shows the mean \ac{utr} transition by time of day, calculated as the weekday mean in four offices.
Generally, \ac{utr} marks the highest values around \DTMtime{08:00:00} to \DTMtime{08:30:00} and around \DTMtime{17:30:00} to \DTMtime{18:00:00}, which are typical occupant arrival and departure times, respectively.
Also, \ac{utr} marks higher values at lunchtime around \DTMtime{12:00:00} to \DTMtime{13:00:00} in shared offices and R\&D Office.
In Shared Office 1 and Shared Office 2, the values change more dynamically than in Admin Office C during office hours.
\ac{utr} marks a higher value around 1.0 on Mondays in R\&D Office, and in Shared Office 2 \ac{utr} marks approximately 0.8 on Mondays.
Except these cases, the \ac{utr} values are generally between 0.2 to 0.5.

\begin{figure}[pos=htbp!]
    \centering
    \includegraphics[width=0.9\linewidth]{\Pic{png}{analysis25_utr30_endpoint_k_scaling_step_rate_full_confirm_rng_k16_trial1_1000_combined_20260416_073500}}
    \caption{Relationship between 30-minute user turnover rate and action-needed window ratio}
    \label{fig:CtrlWindowRatio}
\end{figure}

\begin{figure}[pos=htbp!]
    \centering
    \includegraphics[width=0.9\linewidth]{\Pic{png}{analysis25_utr30_rolling_representative_k_separate_room_panels_30min_full_confirm_rng_k16_trial1_1000_combined_20260416_073500}}
    \caption{Distribution of temperature adjustment magnitudes under dynamic group control}
    \label{fig:UTRexamples}
\end{figure}



% removed
% \SetPicSubDir{ch4-Results/agentic/Result_GEARsimulation}
% As a first part of agentic simulation, realistic one-month simulation was conducted.
% This case is based on actual combination of occupancy logs and comfort models with survey participants.
% An example of daily temperature control and predicted comfort performance both on \ac{sgcm} and \ac{rtgcm} is shown in \autoref{fig:OpeExampls}.

% \autoref{fig:GEARComfPrf}\subref{fig:realComfPro} shows the simulated comfort probability on \ac{sgcm} and \ac{rtgcm} during one-month simulation period. \autoref{fig:GEARComfPrf}\subref{fig:realComfImp} shows the predicted comfort probability improvement from \ac{sgcm} to \ac{rtgcm} during one-month simulation period.

% \begin{figure}[pos=htbp!]
%     \centering
%     \begin{subfigure}{0.97\linewidth}
%         \centering
%         \includegraphics[width=\linewidth]{\Pic{png}{ComfortProbability}}
%         \caption{Simulated comfort probability on \ac{sgcm} and \ac{rtgcm} during one-month simulation period}
%         \label{fig:realComfPro}
%     \end{subfigure}

%     \medskip

%     \begin{subfigure}{0.97\linewidth}
%         \centering
%         \includegraphics[width=\linewidth]{\Pic{png}{ComfortProbabilityImprovement}}
%         \caption{Predicted comfort probability improvement from \ac{sgcm} to \ac{rtgcm} during one-month simulation period}
%         \label{fig:realComfImp}
%     \end{subfigure}
%     \caption{Simulated mean comfort probability of each experiment participant through one-month simulation}
%     \label{fig:GEARComfPrf}
% \end{figure}

% \begin{figure}[pos=htbp!]
%     \centering
%     \includegraphics[width=0.97\linewidth]{\Pic{png}{ComfortAcrossRooms}}
%     \caption{Comfort probabiity ratio across office rooms and whole building on \ac{sgcm} and \ac{rtgcm}}
%     \label{fig:GEARComfPrfRoom}
% \end{figure}
